% Teil II: der durchgerechnete Weg, mit Begründung jedes Schrittes.

\subsection*{Gegeben und gesucht}
Gegeben: $r=\qty{2}{mm}=\qty{2e-3}{m}$,
$m=\qty{0.263}{g}=\qty{2.63e-4}{kg}$,
$\eta_{\text{Honig}}=\qty{10}{Pa\,s}$,
$\rho_f=\qty{1.4}{kg/L}=\qty{1400}{kg/m^3}$, $g=\qty{9.81}{m/s^2}$.
Gesucht: (a) die wirkenden Kräfte, (b) die Zeit $t_*$, bis $v=0.99\,v_\infty$.

\subsection*{Logik in einem Satz}
Die Reibung ist proportional zur Geschwindigkeit ($F_R=6\pi\eta r\,v$); darum ist
die Beschleunigung ``Antrieb minus etwas Ständig proportional zu $v$'' --- eine
lineare DGL mit Endgeschwindigkeit $v_\infty=A/B$ und
Einschwingzeit $t_*=\ln(100)/B$ für $99\,\%$.

\subsection*{Schritt 1 --- Kräfte benennen (Teil a)}
\begin{itemize}
  \item $F_g=mg$, nach unten ($+$):
        $F_g=\qty{0.263e-3}{kg}\cdot\qty{9.81}{m/s^2}\approx\qty{2.58}{mN}$.
  \item $F_A=\rho_f\,V\,g$, nach oben: Auftrieb.
  \item $F_R=6\pi\eta r\,v$, nach oben: Stokes-Reibung, wächst mit $v$.
\end{itemize}
Newton: $m\dot v=F_g-F_A-F_R$. Der Auftrieb schwächt das Gewicht; die Reibung
wächst mit der Geschwindigkeit an, bis der Nettowert auf null zurückgeht.

\subsection*{Schritt 2 --- Volumen und die Konstanten $A$, $B$}
Kugelvolumen:
\[
 V=\frac43\pi r^3=\frac43\pi (2\times10^{-3})^3
   \approx 3.35\times10^{-8}\,\mathrm{m^3}.
\]
Damit die zwei Konstanten der DGL $\dot v = A - Bv$:
\begin{align}
  A&=g\Big(1-\frac{\rho_f V}{m}\Big)
    = \qty{9.81}{m/s^2}\Big(1-\frac{1400\cdot 3.35e-8}{2.63e-4}\Big)\\
    &\approx 9.81\,(1-0.178)\approx \qty{8.06}{m/s^2},\\[2mm]
  B&=\frac{6\pi\eta r}{m}
    =\frac{6\pi\cdot 10\cdot 2\times10^{-3}}{2.63\times10^{-4}}
    \approx 1.43\times10^{3}\,\mathrm{s^{-1}}.
\end{align}

\subsection*{Schritt 3 --- Endgeschwindigkeit}
Für $t\to\infty$ verschwindet $e^{-Bt}$, also
\[
 v_\infty=\frac{A}{B}
   \approx \frac{8.06}{1.43\times10^3}
   \approx \qty{5.6e-3}{m/s}
   =\qty{5.6}{mm/s}.
\]
Bedeutung: Bei etwa $5.6\,\mathrm{mm/s}$ heben sich Antrieb und Reibung exakt
auf --- die Kugel sinkt danach konstant langsam.

\subsection*{Schritt 4 --- Zeit bis $99\,\%$ der Endgeschwindigkeit}
Mit $v(0)=0$ ist $v(t)=v_\infty\big(1-e^{-Bt}\big)$. Gefordert
$v=v_\infty\cdot 0.99$:
\[
 1-e^{-Bt_*}=0.99
 \;\Rightarrow\;
 e^{-Bt_*}=0.01
 \;\Rightarrow\;
 t_*=\frac{\ln(100)}{B}
   \approx \frac{4.605}{1.43\times10^3}
   \approx \qty{3.2}{ms}.
\]
Interpretation: Bei zähem Honig ``einschwingt'' die Kugel in gut
$3\,\mathrm{ms}$ praktisch vollständig auf ihre Endgeschwindigkeit.

\begin{ergebnis}
\begin{equation*}
 F_g\approx\qty{2.58}{mN},\quad
 F_A=\rho_f V g\approx\qty{0.46}{mN},\quad
 F_R=6\pi\eta r\,v\ (\text{wächst mit } v),
\end{equation*}
\begin{equation*}
 v_\infty\approx\qty{5.6}{mm/s},
 \qquad
 t_*\approx\qty{3.2}{ms}.
\end{equation*}
\end{ergebnis}

\begin{kontrolle}
\textbf{Einheiten:} $B=6\pi\eta r/m$:
$\frac{Pa\cdot m}{kg}=\frac{kg\,m\,s}{s^2\,m^2}\,\frac{m}{kg}=\mathrm{s^{-1}}$;
damit ist $t_*=\ln100/B$ in $\mathrm{s}$. \checkmark

\textbf{Querprobe (Grenzfall):} $v_\infty\propto (m-\rho_f V)/\eta$ --- schwerere
Kugel oder dünneres Medium $\Rightarrow$ grössere Endgeschwindigkeit;
$t_*\propto m/(\eta r)$ --- kleinere Kugel oder grössere Viskosität
$\Rightarrow$ schnelleres Einschwingen. Beides passt zur Intuition.

\textbf{Auftriebs-Anteil:} $\rho_f V/m = 1400\cdot 3.35e-8 / 2.63e-4\approx
0.178$ --- der Auftrieb schwächt den Antrieb um $\sim 18\,\%$; ohne ihn wäre
$v_\infty$ entsprechend grösser.

\textbf{Grössenordnung:} mm/s und ms bei zähem Honig --- typisch für eine
Stokes-Bewegung. \checkmark
\end{kontrolle}

\begin{merksatz}
Lineare Reibung erzeugt eine Relaxierung $v\to v_\infty$ mit der
Zeitkonstanten $\tau=1/B=m/(6\pi\eta r)$. Die Zeit bis zum Bruchteil $p$ ist
\[
 t_*=\frac{1}{B}\ln\frac{1}{1-p},
\]
und hängt nur von $B$ und $p$ ab --- nicht von $v_\infty$. Genau das steckt
in der Übungs-Gleichung $x'=A-Bx$ (dort ist $x$ die Geschwindigkeit $v$).
\end{merksatz}
